\section{Related work}

Electric-bus scheduling jointly assigns service to vehicles and tracks battery and charging decisions. Wu et al. study multi-depot scheduling with grid objectives; de Vos et al. model partial charging and station capacity; Parmentier et al. scale route and recharge scheduling by column generation \citep{wu2022grid,devos2024capacity,parmentier2023fleets}. These works establish operational roles for duty construction, energy limits, and charging resources. Here we ask how a complete fleet plan relates to the convex hull of complete plans under a common convex supply-cost model. A duty-column relaxation need not equal this whole-fleet hull, and a mean load in the hull need not be executable as one day's bus schedule.

Public timetable and movement data make this distinction testable without fixing the economics. Sistig et al. generate bus and crew schedules for twenty German transport networks from public GTFS inputs; their versioned dataset includes routes, trips, stops, possible deadheads, and schedules \citep{sistig2025dataset,sistig2025electrification}. We use one Hildenbrand timetable in two restricted depot variants. Source distances, times, and bus parameters are combined with EGG's declared energy, battery, reserve, charger, horizon, and synthetic convex supply-cost assumptions; the dataset supplies no EGG market-supply function. These remain small development studies on one timetable and two depots, with public global bounds open. They establish neither an equal-quality speedup nor a learned-proposal benefit.

Nonconvex electricity-market research explains why marginal prices can fail to support dispatch with indivisible operating decisions. Gribik et al. compare unit-commitment pricing rules and discuss energy uplift when no energy-price vector supports selected quantities \citep{gribikHoganPope2007}; O'Neill et al. construct prices from a mixed-integer model and associated linear program \citep{oneill2005}. Convex-hull pricing studies uniform prices through convexification: Schiro et al. discuss its formulation and implementation, and Hua and Baldick give a convex primal formulation \citep{schiro2016,huaBaldick2017}. Madani et al. compare convex-hull, integer-programming, and European-style pricing with nonconvex demand bids; Andrianesis et al. compute convex-hull prices by Dantzig--Wolfe decomposition and column generation \citep{madani2018pricing,andrianesis2022hull}. Hümbs et al. characterize linear-price support for decentralized problems with integralities and complete linear feasible-set descriptions \citep{huembs2022complete}. These are precedents for convex-hull pricing and price-support questions; we use the same general supporting-price logic for a stated fleet set and supply-cost function. Our own-price best-response test is conditional on that model, not a general market-clearing theorem or settlement design; fleet regret is not an uplift payment under these market rules.

Aggregation results also require their own scale and participant definitions. Starr's classical quasi-equilibrium result gives an aggregation perspective for economies with nonconvex preferences \citep{starr1969}. Our replication construction is proved directly for its specified fleet schedules: the aggregate planning gap tends to zero while whole-operator regret can remain positive, whereas per-bus and individually reserved-pair incentives vanish. A Shapley--Folkman aggregation argument does not by itself prove this exact replication result or force the absolute regret of one composite operator to vanish. The distinction is between aggregating many separately bounded participants and enlarging a single participant whose opportunity set remains indivisible.

Price coordination for flexible EV charging is established. Gan et al. model each vehicle's charging profile as a continuous, rate-bounded vector with fixed energy over an availability window and give a price-update protocol converging for a convex aggregate-load objective \citep{ganTopcuLow2013}. This coordinates charging once availability and energy needs are fixed. Bus schedules additionally assign trips to vehicles, changing charging windows and the joint feasible load set. Convex charging coordination therefore does not imply own-price support for a discrete complete bus schedule.

Fixed-schedule fleet counting has a classical analogue in Dantzig and Fulkerson's minimum-tanker problem \citep{dantzigFulkerson1954}. Such flow reductions count vehicles covering fixed movements, but do not certify battery states, shared-charger feasibility, or price support. We separate minimum-fleet flow diagnostics, charging replay, and convex-hull economic comparisons.
